\documentclass[UTF8,11pt]{ctexart}
\usepackage[a4paper,margin=24mm]{geometry}
\usepackage{amsmath,amssymb,amsthm,mathtools,mathrsfs}
\usepackage[all]{xy}
\usepackage{xurl,hyperref,enumitem}
\usepackage{tikz}
\usetikzlibrary{arrows.meta,cd,decorations.markings,calc,patterns}
\hypersetup{hidelinks,breaklinks=true}
\setlength{\emergencystretch}{3em}
\allowdisplaybreaks
\newtheorem*{theorem}{Theorem}
\newtheorem*{thm}{Theorem}
\newtheorem*{lemma}{Lemma}
\newtheorem*{proposition}{Proposition}
\newtheorem*{prop}{Proposition}
\newtheorem*{corollary}{Corollary}
\newtheorem*{cor}{Corollary}
\newtheorem*{claim}{Claim}
\theoremstyle{definition}
\newtheorem*{definition}{Definition}
\newtheorem*{defn}{Definition}
\newtheorem*{remark}{Remark}
\newtheorem*{example}{Example}
\newcommand{\R}{\mathbb R}
\newcommand{\C}{\mathbb C}
\newcommand{\Z}{\mathbb Z}
\newcommand{\Q}{\mathbb Q}
\newcommand{\N}{\mathbb N}
\newcommand{\F}{\mathbb F}
\newcommand{\D}{\mathbb D}
\newcommand{\bB}{\mathbb B}
\newcommand{\bC}{\mathbb C}
\newcommand{\bR}{\mathbb R}
\newcommand{\bZ}{\mathbb Z}
\newcommand{\bN}{\mathbb N}
\newcommand{\bQ}{\mathbb Q}
\newcommand{\bD}{\mathbb D}
\newcommand{\bF}{\mathbb F}
\newcommand{\CP}{\mathbb{CP}}
\newcommand{\RP}{\mathbb{RP}}
\newcommand{\sO}{\mathscr O}
\newcommand{\sI}{\mathscr I}
\newcommand{\sF}{\mathscr F}
\newcommand{\sG}{\mathscr G}
\newcommand{\sH}{\mathscr H}
\newcommand{\sR}{\mathscr R}
\newcommand{\sM}{\mathscr M}
\newcommand{\sV}{\mathscr V}
\newcommand{\sU}{\mathscr U}
\DeclareMathOperator{\Hom}{Hom}
\DeclareMathOperator{\Ext}{Ext}
\DeclareMathOperator{\Tor}{Tor}
\DeclareMathOperator{\Spec}{Spec}
\DeclareMathOperator{\Proj}{Proj}
\DeclareMathOperator{\supp}{supp}
\DeclareMathOperator{\im}{im}
\DeclareMathOperator{\id}{id}
\DeclareMathOperator{\Tr}{Tr}
\DeclareMathOperator{\tr}{tr}
\DeclareMathOperator{\rank}{rank}
\DeclareMathOperator{\ord}{ord}
\DeclareMathOperator{\dist}{dist}
\DeclareMathOperator{\End}{End}
\DeclareMathOperator{\Aut}{Aut}
\DeclareMathOperator{\Gal}{Gal}
\DeclareMathOperator{\vol}{vol}
\DeclareMathOperator{\Cl}{Cl}
\DeclareMathOperator{\coker}{coker}
\DeclareMathOperator{\codim}{codim}
\DeclareMathOperator{\divergence}{div}
\DeclareMathOperator{\Ric}{Ric}
\DeclareMathOperator{\grad}{grad}
\DeclareMathOperator{\Hess}{Hess}
\newcommand{\abs}[1]{\left\lvert#1\right\rvert}
\newcommand{\sabs}[1]{\lvert#1\rvert}
\newcommand{\babs}[1]{\bigl\lvert#1\bigr\rvert}
\newcommand{\Babs}[1]{\Bigl\lvert#1\Bigr\rvert}
\newcommand{\bbabs}[1]{\biggl\lvert#1\biggr\rvert}
\newcommand{\BBabs}[1]{\Biggl\lvert#1\Biggr\rvert}
\newcommand{\norm}[1]{\left\lVert#1\right\rVert}
\newcommand{\snorm}[1]{\lVert#1\rVert}
\newcommand{\bnorm}[1]{\bigl\lVert#1\bigr\rVert}
\newcommand{\Bnorm}[1]{\Bigl\lVert#1\Bigr\rVert}
\newcommand{\bbnorm}[1]{\biggl\lVert#1\biggr\rVert}
\newcommand{\BBnorm}[1]{\Biggl\lVert#1\Biggr\rVert}
\newcommand{\widebar}[1]{\overline{#1}}
\newcommand{\nicefrac}[2]{\frac{#1}{#2}}
\newcommand{\linnprod}[2]{\langle#1,#2\rangle}
\newcommand{\myindex}[1]{#1}
\newcommand{\glsadd}[1]{}
\newcommand{\avoidbreak}{}
\newcommand{\sourceref}[1]{\textnormal{[原文：\nolinkurl{#1}]}}
\newcommand{\thmref}[1]{\sourceref{#1}}
\newcommand{\lemmaref}[1]{\sourceref{#1}}
\newcommand{\propref}[1]{\sourceref{#1}}
\newcommand{\corref}[1]{\sourceref{#1}}
\newcommand{\defnref}[1]{\sourceref{#1}}
\newcommand{\exampleref}[1]{\sourceref{#1}}
\newcommand{\exerciseref}[1]{\sourceref{#1}}
\newcommand{\figureref}[1]{\sourceref{#1}}
\newcommand{\sectionref}[1]{\sourceref{#1}}
\newcommand{\appendixref}[1]{\sourceref{#1}}
\newcommand{\stacksref}[2]{\href{https://stacks.math.columbia.edu/tag/#1}{#2 [#1]}}


\newcommand{\E}{\mathbb E}
\newcommand{\Prob}{\mathbb P}
\newcommand{\Reff}{\mathcal R_{\mathrm{eff}}}
\newcommand{\eps}{\varepsilon}
\DeclareMathOperator{\diam}{diam}
\DeclareMathOperator{\Var}{Var}
\DeclareMathOperator{\Crit}{Crit}
\newcommand{\dd}{\,\mathrm d}

\begin{document}
\section*{066\quad 一维整环有限分式域扩张中的Krull--Akizuki定理}
\subsection*{来源与计数目标}
The Stacks Project Authors, \emph{The Stacks Project}, Section 10.119，主命题 Lemma 10.119.12，支撑 Lemmas 10.119.9、10.119.11。使用2026-09-28实际访问的在线正文。
\par\url{https://stacks.math.columbia.edu/tag/00P7}
\par\textbf{本文件只计一个问题：}若 $R$ 为一维 Noetherian 整环，$L$ 为其分式域的有限扩张，则每个中间环 $R\subset A\subset L$ 都 Noetherian。
\subsection*{作者命题与人类证明}
\begin{lemma}[10.119.9]
Let $R$ be a domain with fraction field $K$. Let $M$ be an $R$-submodule of $K^{\oplus r}$. Assume $R$ is local Noetherian of dimension $1$. For any nonzero $x\in R$ we have $\operatorname{length}_R(R/xR)<\infty$ and
\[\operatorname{length}_R(M/xM)\leq r\operatorname{length}_R(R/xR).\]
\end{lemma}
\begin{proof}
If $x$ is a unit then the result is true. Hence we may assume $x\in\mathfrak m$ the maximal ideal of $R$. Since $x$ is not zero and $R$ is a domain we have $\dim(R/xR)=0$, and hence $R/xR$ has finite length. Consider $M\subset K^{\oplus r}$ as in the lemma. We may assume that the elements of $M$ generate $K^{\oplus r}$ as a $K$-vector space after replacing $K^{\oplus r}$ by a smaller subspace if necessary.

Suppose first that $M$ is a finite $R$-module. In that case we can clear denominators and assume $M\subset R^{\oplus r}$. Since $M$ generates $K^{\oplus r}$ as a vectors space we see that $R^{\oplus r}/M$ has finite length. In particular there exists an integer $c\geq0$ such that $x^cR^{\oplus r}\subset M$. Note that $M\supset xM\supset x^2M\supset\ldots$ is a sequence of modules with successive quotients each isomorphic to $M/xM$. Hence we see that
\[n\operatorname{length}_R(M/xM)=\operatorname{length}_R(M/x^nM).\]
The same argument for $M=R^{\oplus r}$ shows that
\[n\operatorname{length}_R(R^{\oplus r}/xR^{\oplus r})=\operatorname{length}_R(R^{\oplus r}/x^nR^{\oplus r}).\]
By our choice of $c$ above we see that $x^nM$ is sandwiched between $x^nR^{\oplus r}$ and $x^{n+c}R^{\oplus r}$. This easily gives that
\[\begin{split}
r(n+c)\operatorname{length}_R(R/xR)&\geq n\operatorname{length}_R(M/xM)\\
&\geq r(n-c)\operatorname{length}_R(R/xR).
\end{split}\]
Hence in the finite case we actually get the result of the lemma with equality.

Suppose now that $M$ is not finite. Suppose that the length of $M/xM$ is $\geq k$ for some natural number $k$. Then we can find
\[0\subset N_0\subset N_1\subset N_2\subset\ldots N_k\subset M/xM\]
with $N_i\ne N_{i+1}$ for $i=0,\ldots,k-1$. Choose an element $m_i\in M$ whose congruence class mod $xM$ falls into $N_i$ but not into $N_{i-1}$ for $i=1,\ldots,k$. Consider the finite $R$-module $M'=Rm_1+\ldots+Rm_k\subset M$. Let $N'_i\subset M'/xM'$ be the inverse image of $N_i$. It is clear that $N'_i\ne N'_{i+1}$ by our choice of $m_i$. Hence we see that $\operatorname{length}_R(M'/xM')\geq k$. By the finite case we conclude $k\leq r\operatorname{length}_R(R/xR)$ as desired.
\end{proof}

\begin{lemma}[10.119.11]
Let $R$ be a domain with fraction field $K$. Let $M$ be an $R$-submodule of $K^{\oplus r}$. Assume $R$ is Noetherian of dimension $1$. For any nonzero $x\in R$ we have $\operatorname{length}_R(M/xM)<\infty$.
\end{lemma}
\begin{proof}
Since $R$ has dimension $1$ we see that $x$ is contained in finitely many primes $\mathfrak m_i$, $i=1,\ldots,n$, each maximal. Since $R$ is Noetherian we see that $R/xR$ is Artinian and $R/xR=\prod_{i=1,\ldots,n}(R/xR)_{\mathfrak m_i}$ by Proposition 10.60.7 and Lemma 10.53.6. Hence $M/xM$ similarly decomposes as the product $M/xM=\prod(M/xM)_{\mathfrak m_i}$ of its localizations at the $\mathfrak m_i$.

By Lemma 10.119.9 applied to $M_{\mathfrak m_i}$ over $R_{\mathfrak m_i}$ we see each $M_{\mathfrak m_i}/xM_{\mathfrak m_i}=(M/xM)_{\mathfrak m_i}$ has finite length over $R_{\mathfrak m_i}$. Thus $M/xM$ has finite length over $R$ as the above implies $M/xM$ has a finite filtration by $R$-submodules whose successive quotients are isomorphic to the residue fields $\kappa(\mathfrak m_i)$.
\end{proof}

\begin{lemma}[10.119.12, Krull--Akizuki]
Let $R$ be a domain with fraction field $K$. Let $L/K$ be a finite extension of fields. Assume $R$ is Noetherian and $\dim(R)=1$. In this case any ring $A$ with $R\subset A\subset L$ is Noetherian.
\end{lemma}
\begin{proof}
Let $I\subset A$ be a nonzero ideal. By Lemma 10.30.8 we can find a nonzero element $x\in I\cap R$. Then we get $I/xA\subset A/xA$. By Lemma 10.119.11 the $R$-module $A/xA$ has finite length as an $R$-module. Hence $I/xA$ has finite length as an $R$-module. Hence $I$ is finitely generated as an ideal in $A$.
\end{proof}
\subsection*{整理说明（非作者正文）}
长度估计包括有限模与非有限模两种情形，并取入其从局部到全局的论证；没有将整个关键估计当作现成结论。标准先修为零维 Noetherian 环的 Artinian 分解、模长度的基本性质、局部化，以及代数扩张中非零理想与底整环交非零的 10.30.8（Tag 0H7L）。辅助结果只服务本题，不另计数量。

难度依据：不预设中间环或相应模是有限生成的，要先用有限秩分式域空间中的长度比较控制任意子模，再经有限局部化分解返回理想有限生成。原语言和顺序保留；网页提供数学 TeX，包内另存正文转录，不是原章节文件。本题与正规性判据不同，目标不包含整闭性或 $R_1/S_2$ 判别。
\subsection*{可修改的简短标注（非作者正文）}
$c$：代数扩张中间环的Noetherian性。

$a$：分式域向量空间；模长度；清除分母；有限子模逼近；Artinian环分解；局部化；理想有限生成。

\texttt{cross\_theory\_status = uncertain}。当前证明以有限秩向量空间中的模长度控制返回环的有限生成性质；这项联系在 10.119.9、10.119.11 中明确，但是否作为跨理论样本保留进一步标签判断。
标签为非穷尽、可修改初稿。
\subsection*{许可与修改记录}
原作者及作品见来源栏；来源采用 GNU Free Documentation License 1.2 或后续版本。本题证摘录和附加整理沿用该许可。2026-09-28：增加题号、来源定位、独立排版、整理说明和初稿标注；数学内容的取材范围及排版变化见整理说明。
\end{document}
